%
% ###################################################################
%  Copyright (c) 2026, REPLACE: Your Name
%  Editors: A.D. Rollett & M. De Graef
%  All rights reserved.
%
% Licensed under the Creative Commons CC BY-NC-SA 4.0 License,
% hereafter referred to as the "License"; you may not use this
% document except in compliance with the License. You may obtain
% a copy of the License at
%     https://creativecommons.org/licenses/by-nc-sa/4.0/legalcode
% Unless required by applicable law or agreed to in writing, all
% material distributed under the License is distributed on an
% "AS IS" BASIS, WITHOUT WARRANTIES OR CONDITIONS OF ANY KIND,
% either express or implied. See the License for the specific
% language governing permissions and limitations under the License.
% ###################################################################

% ###################################################################
% AUTHOR SETUP — Edit the lines below
% ###################################################################

% Uncomment the next line ONLY if this is a foundational chapter:
%\corechapter{Yes}

% Your name and affiliation (appears below the chapter title):
\OCchapterauthor{Marc De Graef, Carnegie Mellon University}

% Chapter abbreviation: choose a UNIQUE 6-character uppercase code.
% This is used as a prefix for all labels in this chapter.
% Examples: LINALG, TENSOR, NUMSYS, PROJTS
\renewcommand{\chabbr}{DIFGEO}

% Chapter header image (2480 x 1240 pixels, 300 dpi, RGB format).
% Replace \noheaderimage with your header filename (e.g., MYCHAPheader.pdf).
% The file should be placed in chapter/pdf/.
% Note that \graphicspath is defined in the main.tex file.
%\chapterimage{\graphicspath MYCHAPheader.pdf}
\chapterimage{\noheaderimage}

% ###################################################################
% CHAPTER TITLE — Edit the title below
% ###################################################################
\chapter{The Geometry of Diffraction Processes}\label{\chabbr:ch:DiffractionGeometry}

% Register author information (repeat for each co-author)
\writeauthor{\chabbr:ch:DiffractionGeometry}{The Geometry of Diffraction Processes}{De Graef}{Marc}{Materials Science and Engineering}{Carnegie Mellon University}{mdg@andrew.cmu.edu}{https://www.mse.engineering.cmu.edu/directory/bios/degraef-marc.html}

% ###################################################################
% LEARNING OBJECTIVES
% ###################################################################
% Each chapter begins with Learning Objectives.  Each objective should:
% - Contain an active verb (Learn, Understand, Derive, Apply, etc.)
% - Link to the relevant section using the \chabbr:sec:label format
% - Use the OCBurntOrange color for the reference
\begin{learningobjectives}
{\begin{itemize}
    \item[{\color{OCBurntOrange}\ref{\chabbr:sec:introduction}:}] Understand the basic concepts presented in this chapter
    \item[{\color{OCBurntOrange}\ref{\chabbr:sec:methods}:}] Learn the methods used for analysis
    \item[{\color{OCBurntOrange}\ref{\chabbr:sec:examples}:}] Apply the concepts to worked examples
\end{itemize}}
\end{learningobjectives}

% ###################################################################
% CHAPTER CONTENT STARTS HERE
% ###################################################################

%Chapter structure
% - Introduction
% - Particle Waves
% - Wave Interference
% - Bragg's Law in Direct Space
% - Bragg's Law in Reciprocal Space
% -    Ewald sphere construction
% -    Unit vector form of Bragg's Law
% - Important Diffraction Geometries
% 
% 


\section{Introduction}\label{\chabbr:sec:introduction}

In this chapter we will describe the geometry of the diffraction process in a way that includes the most commonly used types of incident beams: x-rays, electrons, and neutrons.  We begin by introducing the concept of the \indexit{wave vector}, which relates the wave length of the incident beam to the reciprocal space description of the crystal structure. This involves the \indexit{particle-wave duality} which is fundamental to the field of quantum mechanics. Then we describe how the three primary beam types are generated.  We build up a description of x-ray scattering by first looking at scattering of an x-ray photon by a single electron, and then by an atom



\section{Particle and Waves}\label{\chabbr:sec:particlewaves}


\section{Wave Interference}\label{\chabbr:sec:interference}


\section{Bragg's Law in Direct Space}\label{\chabbr:sec:BraggDirect}


\section{Bragg's Law in Reciprocal Space}\label{\chabbr:sec:BraggReciprocal}


\section{Important Diffraction Geometries}\label{\chabbr:sec:DiffractionGeometries}






% ###################################################################
% END OF CHAPTER
% ###################################################################
